<<<P000692>%
Это ошибка времени компиляции, если формальный параметр имеет модификатор <<P000823>
Модификатор <<P000824>.
Ошибка времени компиляции <<<P000825> происходит как
Первый токен <<<P000826>{fieldFormalParameter}.

<<<P000693>%
Ошибка времени компиляции, если параметр
<<<P000827><<Формальное поле Параметр возникает как параметр функции
который не является неперенаправляющим генеративным конструктором.

<<<P000828>{%>
A <<<P000829>{fieldFormalParameter} объявляет инициализацию формальной,
который описан в другом месте
(<<P000457>)%
?

<<<P000694>%
Можно включить модификатор \COVARIANT{}
в некоторых формах декларирования параметров.
Эффект от этого описан в отдельном разделе.
(P000458).

<<<P000831><%
Обратите внимание, что также используется не терминальный <<<P000832>{нормальный формальный параметр}
в правилах грамматики для факультативных параметров,
Это означает, что такие параметры также могут быть ковариантными.
?

<<<P000695>%
Это ошибка времени компиляции, если модификатор <<P000833> происходит
параметр функции, который не является
Способ экземпляра, установщик экземпляра или оператор экземпляра.


<<P000834> Факультативные формы
<<P000790>

<<<P000696>%
Дополнительные параметры могут быть указаны и снабжены значениями по умолчанию.

<<P000431>
<defaultFormalParameter> ::= <normalFormalParameter> ('=' <expression>)

<defaultNamedParameter> ::= <<<P000835>
<metadata> <<P000836>? <normalFormalParameterNoMetadata>
\gnewline{}('=' | ':') <выражение>)?
<<P000444>

Форма <<P000838>{<нормальный Формальный Параметр>:> <выражение>}
Это эквивалентно форме
<<P000839>{<normalFormalParameter> '=' <expression>}.
Колонно-синтакс включен только для обратной совместимости.
Он обесценивается и удаляется в
Более поздняя версия спецификации языка.

<<<P000697>%
Это ошибка времени компиляции, если значение по умолчанию необязательного параметра
Не является постоянным выражением (P000459).
Если по умолчанию явно не указан необязательный параметр
Предусмотрен неявный дефолт <<<P000840>{}.

<<<P000698>%
Ошибка времени компиляции, если имя названного опционального параметра
Начинается с символа «_».

<<<P000841><%
Необходимость такого ограничения является прямым следствием
Дело в том, что именование и приватность не являются ортогональными.
Если мы позволим названным параметрам начаться с подчеркивания,
Они считаются частными и недоступными для
Абоненты из-за пределов библиотеки, где она была определена.
Если метод вне библиотеки переопределяется
метод с частным факультативным названием,
Это не будет подтипом оригинального метода.
Статическая шашка, конечно, будет отмечать такие ситуации.
Но последствием будет то, что добавление частного, названного формальным, сломается.
клиенты за пределами библиотеки таким образом, что они не могли легко исправить.
?


<<<P000842>{Ковариантные параметры}
<<P000791>

<<<P000699>%
Дарт допускает формальные параметры примеров,
В том числе сеттеры и операторы,
Заявлено <<P000843>.

<<<P000844>{%>
Синтаксис для этого указан в предыдущем разделе.
(<<P000460>)%
?

<<<P000700>%
Ошибка времени компиляции, если модификатор <<P000845> происходит
в декларации формального параметра функции
который не является методом экземпляра, множителем экземпляров или оператором.

<<<P000846>{%>
Как указано ниже, параметр также может быть ковариантным по другим причинам.
Общий эффект наличия ковариантного параметра <<<P000008>
в подписи данного метода <<<P000009>
- позволяет ковариантно перекрывать тип <<<P000010>,
Это означает, что тип, требуемый во время выполнения для данного фактического аргумента.
может быть подходящим подтипом типа, который известен во время компиляции.
на сайте call.%
?

<<<P000847>{%>
Этот механизм позволяет разработчикам явно запрашивать
Гарантия времени компиляции, которая в противном случае поддерживается(а именно: что фактический аргумент, статический тип которого удовлетворяет требованию)
Это также будет сделано во время выполнения
Заменяются динамическими проверками.
В обмен на принятие этих динамических проверок типа,
Разработчики могут использовать ковариантные параметры для выражения дизайна программного обеспечения.
где динамические проверки типа известны (или, по крайней мере, доверены) для успеха;
на основании рассуждений о том, что анализ статического типа не захватывает.%
?

<<<P000701>%
<<<P000808>%
<<<P000011> Сигнатура метода с формальными параметрами типа
<<P000848> X}{1}{}
Формальные позиционные параметры <<<P000849>{p}{1}{n}
<<<P000809>%
именованные формальные параметры <<<P000850>{q}{1}{k}.
<<<P000810>%
<<<P000012> быть подписью метода с формальными параметрами типа
<<P000851> XX \!}{1}{}
Формальные позиционные параметры <<<P000852><p> '\!}{1}{n},
<<<P000811>%
именованные формальные параметры <<<<P000853>> '\!}{1}{k'}.
%
Предположим, что <<<P000013> и <<<P000014>;
Мы говорим: <<<P000015> является параметром в <<<P000016> это
<<<P000854>{соответствует параметру}
Формальный параметр <<P000017> в <<P000018>.
Предположим, что <<<P000019> и <<<P000020>;
Мы говорим, что <<<P000021> является параметром в <<<P000022> это
<<<P000855>{соответствует} формальному параметру
<<<P000023> в <<<P000024>, если <<<P000025>.
%
Точно так же мы говорим, что формальный параметр типа
<<<P000026> <<<P000027>
<<<P000856>{соответствует} формальному параметру типа
<<<P000028> от <<<P000029>, для всех <<<P000030>.

<<<P000857>{%>
Сюда относится случай, когда <<<P000031> соответственно <<P000032> иметь
необязательные позиционные параметры,
в каком случае <<P000033> соответственно <<P000034> Должен держать,
Но мы можем иметь <<<P000035>.
Случай, когда числа формальных параметров типа отличаются, не имеет значения.
?

% Быть ковариантным — свойство параметра интерфейса класса;
% это означает, что мы говорим только о исходном ключевом слове <<<P000858>
% и класс, содержащий соответствующую декларацию, когда мы обнаруживаем
% первый раз, когда данный параметр является ковариантным. С этого момента и на
% оно «переносится» по ссылкам подтипа, связанным с интерфейсами классов;
% так, что мы можем получить его индуктивно из непрямого суперинтерфейса
% путем проверки того, имеют ли прямые суперинтерфейсы сигнатуру метода
% с соответствующим именем и соответствующим параметром, а затем проверка
% от этого параметра. Такой же подход применим и к ковариантам по классам.

<<<P000702>%
<<<P000812>%
<<<P000036> Это класс, который объявляет
Метод <<<P000037> который имеет
параметр <<P000038> заявление которого имеет модификатор <<P000859>;
В этом случае мы говорим, что параметр <<<P000039> это
<<<P000860>{covariant-by-declaration}{parameter!covariant-by-declaration}.
%
В этом случае интерфейс <<<P000040> имеет сигнатуру метода <<<P000041>,
и эта подпись имеет параметр <<<P000042>;
Мы также говорим, что параметр <<<P000043> В этом методе подпись
<<<P000861>{covariant-by-declaration}
%
Наконец, параметр <<<P000044> сигнатуры метода <<<P000045>
интерфейс класса <<<P000046> это
<<<P000862>{ковариант по декларации}
если прямой суперинтерфейс <<<P000047>
имеет доступную сигнатуру метода <<<P000048> с таким же названием, как <<<P000049>,
который имеет параметр <<<P000050> что соответствует <<<P000051>,
<<<P000052> является ковариантным по декларированию.

<<<P000703>%
Предположим, что <<<P000813><<<P000053> Является общим классом
с официальными декларациями параметров типа
<<P000554>,
<<<P000814>%
<<<P000055> быть декларацией метода экземпляра в <<<P000056>
(который может быть методом, установщиком или оператором),
<<<P000057> быть параметром, объявленным <<<P000058>, и
Пусть <<<P000059> будет объявленным типом <<<P000060>.
%
Параметр <<<P000061> это<<<P000863>{ковариант по классу}{параметр!ковариант по классу}
Если для любого <<<P000062>,
<<P000063> происходит в ковариантном или инвариантном положении в <<<P000064>.
%
В этом случае интерфейс <<<P000065> также имеет сигнатуру метода <<<P000066>,
и эта подпись имеет параметр <<<P000067>;
Мы также говорим, что параметр <<<P000068> В этом методе подпись
<<<P000864>{ковариант по классу}.
Наконец, параметр <<<P000069> сигнатура метода <<<P000070>
интерфейс класса <<<P000071> это
<<<P000865>{ковариант по классу}
если прямой суперинтерфейс <<<P000072>
имеет доступную сигнатуру метода <<<P000073> с таким же названием, как <<<P000074>,
который имеет параметр <<P000075> что соответствует <<<P000076>,
<<<P000077> является ковариантным по классу.

<<<P000704>%
Формальный параметр <<<P000078> это
<<<P000866>{ковариант}{параметр!ковариант}
Если <<<P000079> является ковариантным по декларированию, или <<P000080> является ковариантным по классу.

<<<P000867>{%>
Параметр может быть одновременно
Ковариант по классу и ковариант по классу.
Обратите внимание, что параметр может быть
ковариант по классу или ковариант по классу
на основании заявления в любом прямом или косвенном суперинтерфейсе,
В том числе любой суперкласс:
Приведенные выше определения распространяют эти свойства.
Интерфейс от каждого из его прямых суперинтерфейсов
Но они в свою очередь получат
свойство от их прямых суперинтерфейсов,
и так далее.%
?


<<P000868> Тип функции.
<<P000792>

<<<P000705>%
Этот раздел определяет статический тип, который приписывается
функция, обозначенная декларацией функции,
и динамический тип соответствующего функционального объекта.

<<<P000706>%
В данной спецификации,
обозначение, используемое для обозначения типа функции,
То есть, <<<P000536>,
следует синтаксису языка,
За исключением того, что <<<P000869>>
<<P000870>.
Это означает, что каждый тип функции имеет одну из форм.
<<P000871> T_0
<<P000872> T_0

<<P000873>
где <<P000081> является возвратным типом,
<<P000082> Формальные параметры типа с границами <<P000083>, <<P000084>,
<<P000085> Формальные типы параметров для <<<P000086>,
<<P000087> Названия параметров для <<<P000088>.
Типы негенерических функций охватываются случаем <<<P000089>,
где перечень типовых параметров
\code{\ldots{}>>
В целом опущен.
%
Аналогичным образом факультативные скобки <<<P000555> и <<P000876>>>\{\} опущены
когда нет дополнительных параметров.

% Мы обещаем, что обе формы всегда получают одинаковую обработку для k=0.
<<<P000877>{%>
Обе формы с опционами охватывают типы функций без опционов.
когда <<<P000090>
Каждое правило данной спецификации таково, что
Любая из двух форм может использоваться без двусмысленности.
Определить лечение типов функций без опций.%
?

<<<P000707>%
Если декларация функции не объявляет тип возврата явно,
Тип возврата: \DYNAMIC{} (<<P000461>),
Если это не конструктор,
В этом случае он не считается возвратным,
или является сеттером или оператором <<<P000556>,
В этом случае тип возврата составляет <<P000879>.

<<<P000708>%
Декларация функции может объявлять формальные параметры типа.
Тип функции включает имена параметров типа.
и для каждого параметра типа верхней границы,
который считается встроенным классом <<<P000557>
если не указана граница.
При последовательном переименовании параметров типа
может сделать два типа функций идентичными;
Они считаются одного и того же типа.

<<<P000880>{%>
Удобно включать формальные имена параметров типа в типы функций.
Они необходимы для того, чтобы выразить такие вещи, как отношения между
различные параметры типа, F-связи и типы формальных параметров.Однако мы не хотим различать два типа функций.
если они имеют одинаковую структуру и отличаются только выбором названий.
Эта обработка имен также известна как альфа-эквивалентность.
?

<<<P000709>%%
В следующих трех пунктах,
Если число <<P000881>{s} формальных параметров типа равно нулю, то
Перечень параметров типа в типе функции опущен.

<<<P000710>%
<<<P000091> Быть функцией с
Типовые параметры <<P000882>,
Требуемые формальные типы параметров <<P000883> T}{1}{n}
Тип возврата <<P000092>,
без дополнительных параметров.
Статический тип <<<P000093> это
<<P000884>{T_0}.

<<<P000711>%
<<<P000094> Быть функцией с
Типовые параметры <<P000885>,
Требуемые формальные типы параметров <<P000886> T}{1}{n}
Тип возврата <<P000095>
и позиционные факультативные типы параметров <<<P000887>{T}{n+1}{n+k}.
Тогда статический тип <<<P000096> это
<<<P000888>{T_0}.

<<<P000712>%
<<<P000097> Быть функцией с
Типовые параметры <<P000889>,
Требуемые формальные типы параметров <<P000890> T}{1}{n}
Тип возврата <<P000098>,
именованные параметры <<<P000891>{T}{x}{n+1}{n+k}
где <<<P000099> <<P000100> может иметь или не иметь значение по умолчанию.
Статический тип <<<P000101> это
<<<P000892>{T_0}.

<<<P000713>%
<<<P000102> Статический тип функциональной декларации <<<P000103>.
<<<P000104> быть типом времени выполнения объекта функции <<<P000105> полученный
функция клозурита
(<<P000462>)
или пример метода клозурирования
(<<P000463>)
Применительно к <<<P000106>,
<<<P000107>> быть фактическим типом, соответствующим <<<P000108>
В случае, когда <<<P000109> был создан
(<<P000464>).
<<<<<<<<<<<<<<<<<<<<<<<<<<<<<<<<<<<<<<<<<<<<<<<<<<<<<<<<<<<<<<<<<<<<<<<<<<<<<<<<<<<<<<<<<<<<<<<<<<<<<<<<<<<<<<<<<<<<<<<<<<<<<<<<<<<<<<<<<<<<<<<<<<<<<<<<<<<<<<<<<<<<<<<<<<<<<<<<<<<<<<<<<<<<<<<<<<<<<<<<<<<<<<<<<<<<<<<<<<<<<<<<<<<<<<<<<<<<<<<<<<<<<<<<<<<<<<<<
<<P000110> Может содержать переменные свободного типа, но <<<P000111> Содержит фактические значения.%
?
После этого следует придерживаться следующего:
<<P000112> - класс, реализующий встроенный класс <<<P000894>;
<<P000113> является подтипом <<<P000114>;
<<P000115> не является подтипом какого-либо функционального типа, который является соответствующим подтипом <<<P000116>.
<<<P000895>{%>
Если бы мы пропустили последнее требование,
<<<P000896> <<<P000897>< int\,\FUNCTION[int]]
можно оценить по <<P000899> с заявлением
\code{<<P000901>>>{} f()<<P001276>>>\{\}},
Что, очевидно, не является целью.%
?

<<<P000902>{%>
Реализация зависит от выбора
соответствующее представление для функциональных объектов.
Например, считать, что
Функционный объект, полученный посредством извлечения свойств
рассматривает равенство иначе, чем другие функциональные объекты;
Скорее всего, это другой класс.
Реализации могут также использовать различные классы для функциональных объектов.
В зависимости от типа и типа.
Арити может косвенно зависеть от того, является ли функция
метод экземпляра (с неявным параметром приемника) или нет.
Вариации многообразны и, например,
Нельзя предположить, что два различных объекта функции
обязательно будет иметь тот же тип времени выполнения.%
?


<<P000903> Внешние функции.
<<P000793>

<<<P000714>%
<<<P000904>{внешняя функция}{функция!внешняя функция}
является функцией, орган которой предоставляется отдельно от его декларации.
Внешняя функция может быть
Функция верхнего уровня (<<P000465>),
метод (<<<P000466>, <<P000467>,
геттер (<<P000468>),
Сеттер (<<P000469>),
Неперенаправляемый конструктор
(<<P000470>, <<P000471>.
Внешние функции вводятся через встроенный идентификатор. <<<P000905><
(<<P000472>)
Затем следует подпись функции.

<<<P000906><<<<<<<<<<<<<<<<<<<<<<<<<<<<<<<<<<<<<<<<<<<<<<<<<<<<<<<<<<<<<<<<<<<<<<<<<<<<<<<<<<<<<<<<<<<<<<<<<<<<<<<<<<<<<<<<<<<<<<<<<<<<<<<<<<<<<<<<<<<<<<<<<<<<<<<<<<<<<<<<<<<<<<<<<<<<<<<<<<<<<<<<<<<<<<<<<<<<<<<<<<<<<<<<<<<<<<<<<<<<<<<<<<<<<<<<<<<<<<<<<<<<<<<<
Внешние функции позволяют вводить информацию о типе кода
что не известно компилятору Dart.%
?

<<<P000907>{%>
Примерами внешних функций могут быть внешние функции.
(определяется в C, Javascript и т.д.),
примитивы реализации (как определено системой времени выполнения Dart),
код, который генерируется динамическиИнтерфейс которого статически известен.
Однако абстрактный метод отличается от внешней функции.
как имеет <<<P000908>{no} тело.%
?

<<<P000715>%
Внешняя функция связана с телом посредством
Конкретный механизм реализации.
Попытка вызвать внешнюю функцию
которая не связана со своим телом.
Выбрасывает <<<P000558> или какой-либо его подкласс.

<<<P000716>%
Может быть повышена конкретная ошибка компиляции
в объявлении функции <<<P000909>{}.

<<<P000910><<%
Такие ошибки указывают на то, что
Каждый вызов этой функции будет бросать, например,
потому что известно, что он не будет связан с телом.
?

<<<P000717>%
Настоящий синтаксис приведен в
разделы <<P000473> и <<P000474> ниже.


<<P000911> Классы
<<P000794>

<<<P000718>%
<<<P000537> определяет форму и поведение набора объектов, которые являются его
\IndexCustom{instances}{instance}.
Классы могут быть определены классовыми декларациями, как описано ниже.
или через приложения для смешивания (<<<P000475>).

<<P000432>
<classDeclaration> ::=
<<P000913> <<<P000914 >> <typeIdentifier> <typeParameters>
<<<P000915>{} <суперкласс>? <Интерфейсы>.
\gnewline{} '{' (<metadata> <classMemberDeclaration>)* '} "
<<P000917> <<P000918> \CLASS{{}<mixinApplication Класс

<typeNotVoidList> ::= <typeNotVoid> (',' <typeNotVoid>)*

<classMemberDeclaration> ::= <declaration>; "
<<P000920> <методическая подпись> <Функциональное тело>

<methodSignature> ::= <constructorSignature> <initializers>
<<<P000921> <factoryConstructorSignature>
<<P000922> <<P000923> <Функциональная подпись>
<<P000924> <<P000925> <getterSignature>
<<P000926> <<P000927> <setter подпись>
<<<P000928> <подпись оператора>

<декларация> ::= \EXTERNAL{{}<factoryConstructorSignature>
<<P000930> \EXTERNAL{{}<constantConstructorSignature>
<<P000932> <<<P000933>{}<ConstructorSignature>
<<P000934> (<<P000935>{) <<P000936>? <getterSignature>
<<P000937> (\EXTERNAL) <<P000939>? <setter подпись>
<<P000940> (\EXTERNAL) <<P000942>? <Функциональная подпись>
<<P000943> <<<P000944>? <operatorSignature>
<<P000945> <<<P000946>< <<<P000947><type>? <staticFinalDeclarationList>
<<P000948> \STATIC{} <<<P000950><type>? <staticFinalDeclarationList>
<<P000951> <<<P000952>< <<<P000953>< <<<P000954>{} <type>? <инициализированный ИдентификаторСписок
<<P000955> <<<P000956> <<P000957> <varOrType> <инициализированный ИдентификаторСписок
<<P000958> <<<P000959>< <<<P000960>< <<<P000961><type>? <identifierList>
<<P000962> <<<P000963>< <<P000964> <varOrType> <инициализированный ИдентификаторСписок
<<P000965> <<P000966> <<<P000967 >>{} <type>? <инициализированный ИдентификаторСписок
<<P000968> <<P000969> <varOrType> <инициализированный ИдентификаторСписок
<<<P000970> <redirectoryConstructorSignature>
<<P000971> <constantConstructorSignature> (<redirection> | <initializers>)?
<<P000972> <constructorSignature> (<redirection> | <initializers>)?

<staticFinalDeclarationList>:= <<<P000973><
<staticFinalDeclaration> (<staticFinalDeclaration>) *

<staticFinalDeclaration> ::= <identifier> <expression>
<<P000445>

<<<P000719>%
Можно включить модификатор \COVARIANT{}
в некоторых формах декларации.
Эффект от этого описан в другом месте.
(P000476).

<<<P000720>%
Класс имеет конструкторов, экземпляров и статических членов.
<<<P000975>{члены организации}{члены организации!) класса
являются его методы экземпляра, геттеры, множители и переменные экземпляра.
<<<P000976>{статические члены}{члены!статические} класса
Это его статические методы, геттеры, сеттеры и переменные.
<<<P000977>{члены}{члены} класса
являются его статическими и экземплярными членами.

<<P000721>%
Классовая декларация вводит две сферы:

<<P000433>
<<P000978>A <<<P000979>{type-parameter scope}{scope!type parameter},
который является пустым, если класс не является общим (<<P000477>).
Включающая сфера охвата типового параметра декларации класса является
библиотечный объем текущей библиотеки.
<<P000980>
A <<<P000981>{область тела}{сфера действия} Классное тело.
Включающая сфера охвата органа сфера охвата классовой декларации является
Сфера охвата параметра типа декларации класса.
<<P000446>

<<<P000722>%
текущая сфера действия декларации члена инстанции,
Статическая декларация,
или декларация застройщика является
Область применения класса, в котором он объявлен.

<<<P000723>%
Нынешний случай
<<<P000982>{(и, следовательно, его члены)}
Доступ можно получить только в определенных местах в классе:
Мы говорим, что местоположение <<<P000117>
<<P000983>>>{имеет доступ к <<P000984>>>}{имеет доступ к этому@имеет доступ к <<P000985>>>}
f <<<P000118> находится внутри тела декларации
элемент экземпляра или генеративный конструктор,
или в начальном выражении <<<P000986> Пример переменной декларации.

<<<P000987>{%>
Обратите внимание, что начальное выражение для не-<<P000988>>>{} переменная переменная
не имеет доступа к <<P000989>,
и ни одна из частей декларации с пометкой <<P000990>.%
?

<<<P000724>%
Каждый класс имеет один суперкласс.
кроме классов <<P000559> Который не имеет суперкласса.
Класс может реализовывать несколько интерфейсов.
Объявив их в своем имплементационном пункте (<<P000478>).

<<<P000725>%
<<<P000991>{абстрактная декларация класса}{классовая декларация}! абстрактный
Классовая декларация, которая явно декларируется
<<<P000992>{} Модификатор.
A <<<P000993>{конкретная декларация класса}{классовая декларация}! конкретный
Классовая декларация, которая не является абстрактной.
<<<P000994>{абстрактный класс}{класс!абстракт} - это класс
заявление которого является абстрактным, и
а <<<P000995>{конкретный класс}{класс!конкретный} - это класс
Чья декларация конкретна.

<<<P000996>{%>
Мы хотим различного поведения для конкретных классов и абстрактных классов.
Если <<P000119> Он должен быть абстрактным,
Мы хотим, чтобы статическая шашка предупреждала о любой попытке инстанцировать <<P000120>,
Мы не хотим, чтобы чекеры жаловались на нереализованные методы в P000121.
В отличие от этого, если <<<P000122> Он должен быть конкретным,
Проверяющий должен предупредить обо всех нереализованных методах,
Позволяет клиентам делать это свободно.%
?

<<<P000997>{%>
Интерфейс класса <<<P000123> это
неявный интерфейс, который объявляет подписи членов экземпляра
которые соответствуют членам экземпляра, заявленным <<<P000124>,
и чьи прямые суперинтерфейсы
Прямые суперинтерфейсы <<P000125>
(<<P000479>, <<<P000480>>%
?

<<P000726>%
Когда имя класса появляется как тип,
Это имя обозначает интерфейс класса.

<<<P000727>%
Это ошибка времени компиляции, если класс называется <<<P000126> объявлять
член с базовым именем (<<P000481>) <<P000127>.
Если общий класс называется <<<P000128> объявляет переменную типа <<<P000129>,
Это ошибка времени компиляции
Если <<P000130> равно <<<P000131>,
или если <<P000132> имеет член, базовое имя которого <<<P000133>,
или если <<P000134> Имеет конструктор под названием <<<P000560>.

<<<P000998>{%>
Вот простые примеры, иллюстрирующие разницу между
"имеет члена" и "объявляет его членом".
Например, <<P000561> \IndexCustom{declares}{declares member}
один член, названный <<P000562>,
Но это <<<P001000>{has}{has Member} два таких члена.
Правила наследования определяют, какие члены класса имеют.
?

<<P000000>


<<<P001001> Полностью реализовать интерфейс
<<P000795>

% Обратите внимание, что правила здесь и в <<<P000482> Они пересекаются, но они
% необходимое: Этот раздел посвящен конкретным методам, включая
% унаследованных и <<P000483> Заинтересован в случае% членов, заявленных в <<<P000137>, включая как конкретные, так и абстрактные.

<<<P000728>%
% Использование "конкретного члена" ниже может показаться избыточным, поскольку класс
% не наследует абстрактных членов от своего суперкласса.
% подчеркивает тот факт, что даже при абстрактной декларации <<P000138> это
%, заявленных в <<<P000139>, <<P000140> Не имеет значения <<<P000141> которых здесь будет достаточно.
Конкретный класс должен полностью реализовать свой интерфейс.
<<<P000815>%
<<<P000142> быть конкретным классом, объявленным в библиотеке <<<P000143>, с интерфейсом <<<P000144>.
Предположим, что <<P000145> имеет подпись участника <<<P000146> Доступен по адресу <<<P000147>.
Ошибка времени компиляции <<<P000148> не иметь
конкретный член с тем же названием, что и <<P000149> Доступен по адресу <<<P000150>,
Если только <<P000151> имеет нетривиальный <<<P000563>
(P000484).

<<P000729>%
Каждый конкретный член должен иметь соответствующую подпись:
Предположим, что <<P000152> имеет конкретный член с
Оригинальное название: <<P000153> Доступен по адресу <<<P000154>,
Пусть <<<P001002>{m''} будет его подписью.
<<<P001003>{%>
Конкретный член может быть объявлен в <<<P000155> или унаследованы от суперкласса.
?
Пусть <<<P001004>{m'} является подписью участника, полученной от <<<P000156>
Добавляя, если не присутствует, модификатор <<<P001005>
(<<P000485>)
для каждого параметра <<<P000157> <<<P000158> где
Соответствующий параметр в <<<P000159> Имеет модификатор <<<P001006>.
Ошибка времени компиляции <<<P000160> не является правильным переопределением <<<P000161>
(<<P000486>),
Если этот конкретный элемент не является <<<P000564> экспедитор
(P000487).

<<<P001007>{%>
Рассмотрим конкретный класс <<<P000565>,
Предположим, что <<P000566> Объявлять или наследовать
Реализация с таким же названием
для каждой подписи в интерфейсе.
Это все еще ошибка, если одна или несколько из этих реализаций
имеет такие параметры или типы, которые не удовлетворяют
подпись соответствующего члена в интерфейсе.
Для этой проверки не хватает <<<P001008 >> Модификаторы добавляются неявно
к подписи наследуемого члена (так мы получаем <<<P000162> от <<<P000163>).
При добавлении модификатора <<<P001009>{} к одному или нескольким параметрам
(что произойдет только тогда, когда конкретный член наследуется),
реализация может выбрать неявно индуцировать метод пересылки
с такой же подписью, как <<<P000164>,
для выполнения требуемой динамической проверки типа,
Используйте унаследованный метод.%
?

<<P000730>%
Это конкретный выбор реализации, независимо от того, является ли он или нет.
используется неявно индуцированный метод пересылки
когда модификатор <<<P001010>{} добавляется
один или несколько параметров в <<<P000165>.

<<<P001011>{%>
Это верно, несмотря на то, что
Такие методы пересылки можно наблюдать.
Например, мы можем сравнить
время выполнения типа отрыва метода от
Получатель типа <<P000567>
к типу времени выполнения отрыва суперметода от
местоположение в теле <<<P000568>.%
?

<<<P000731>%
с использованием или без использования метода,
Подпись участника в интерфейсе <<<P000166> является <<<P000167>.

<<<P001012>{%>
Способ пересылки не изменяет интерфейс <<<P000569>,
Это деталь реализации.
В частности, это касается даже тех случаев, когда
четкое заявление о методе пересылки имело бы
Изменился интерфейс <<<P000570>, потому что $m'$ Это подтип <<<P000169>.

Когда класс имеет нетривиальный <<P000571>,
класс может оставить некоторых членов нереализованными,
и классу разрешается иметь <<<P000572> экспедитор
который не удовлетворяет интерфейсу класса
(в этом случае он будет отменен)
Еще один экспедитор P000573.

Вот пример: %
?

<<<P000001>

<<P000732>%
Параметры, которые являются ковариантными по декларированиюдолжны также удовлетворять следующим ограничениям:
Предположим, что параметр <<<P000170> <<<P000171> Имеет модификатор <<<P001013>.
Допустим, что прямой или косвенный суперинтерфейс <<P000172> иметь
Подпись метода <<<P000173> с таким же названием, как <<P000174> Доступен по адресу <<<P000175>,
<<<P000176> имеет параметр <<<P000177> Это соответствует P000178.
В этой ситуации происходит ошибка компиляции времени.
Если тип <<<P000179> Это не подтип и не супертип типа <<<P000180>.

<<<P001014>{%>
Это гарантирует, что унаследованный метод удовлетворяет тем же ограничениям.
для каждого формального параметра, который является ковариантным по декларированию
как ограничение, указанное для декларации в <<<P000181>
(<<P000488>) %
?


<<<P001015> Методы инстанций
<<P000796>

<<P000733>%
\IndexCustom{Методические методы}{метод!инстанция}
Функции (<<P000489>)
декларация которых непосредственно содержится в классовой декларации
И это не объявлено <<<P001017>.
<<<P000538> <<<P000182> есть
Методы экземпляра, объявленные <<<P000183>
и методы экземпляра, унаследованные <<<P000184> Из своего суперкласса
(<<P000490>).

<<<P000734>%
<<<P000816>%
Рассмотрим класс <<P000185>
и заявление члена экземпляра <<P000186> в $C$, с подписью члена $m$
(<<P000491>).
Ошибка времени компиляции <<<P000189> отменяет декларацию
% Обратите внимание, что <<<P000190> является доступной в силу определения "переопределений".
с подписью члена <<P000191>
Прямой суперинтерфейс <<P000192>
(<<P000492>),
Если <<<P000193> не является правильным членом, переопределяющим <<<P000194>
(P000493).

<<<P001018>{%>
Это не единственный вид конфликта, который может существовать:
Заявление члена суда <<<P000195> может противоречить другому заявлению <<P000196>,
Даже в тех случаях, когда они не имеют одинаковых названий.
Или это не одно и то же заявление.
Например, <<P000197> Это может быть пример Геттера и <<<P000198> статический сеттер
(<<P000494>) %
?

<<P000735>%
Для каждого параметра $p$ $m$, где присутствует \COVARIANT{},
Ошибка времени компиляции при наличии
прямой или косвенный суперинтерфейс <<<P000201> который имеет
Подпись доступного метода <<<P000202> с таким же названием, как <<<P000203>,
<<<P000204>> имеет параметр <<<P000205> что соответствует <<<P000206>
(<<P000495>),
Если только тип <<P000207> Это подтип или супертип типа <<<P000208>.

<<<P001020>{%>
Это означает, что
параметр, который является ковариантным по декларированию, может иметь тип
который является супертипом или подтипом типа
соответствующий параметр в суперинтерфейсе,
Но эти два типа не могут быть несвязанными.
Обратите внимание, что это требование должно быть выполнено
для каждого прямого или косвенного суперинтерфейса отдельно,
Потому что эти отношения не транзитивны.
?

<<<P001021>{%>
Суперинтерфейс может быть статически известным типом приемника.
Это означает, что мы расслабляем потенциальные отношения типизации.
между статически известным типом параметра и
Тип, который действительно требуется во время выполнения
отношения подтипа или супертипа,
Вместо строгих отношений супертипа
Это относится к параметру, который не является ковариантным.
Следует отметить, что статически он не известен.
на месте вызова, является ли любой данный параметр ковариантным,
Поскольку ковариация может быть введена в
надлежащий подтип статически известного типа приемника.
Мы решили отдать приоритет гибкости, а не безопасности здесь.
Потому что весь смысл ковариантных параметров в том, что разработчики
Вы можете сделать выбор, чтобы увеличить гибкость.
в компромиссе, где некоторая безопасность статического типа теряется.
?


<<<P001022> Операторы
<<P000797>

<<<P000736>%
<<<P001023>{Операторы}{операторы} - это методы экземпляров со специальными именами,За исключением оператора \lit{[]] который является экземпляром Геттера
и оператор <<<<P001025>{[]=}, который является набором экземпляров.

<<P000434>
<оператор подпись> ::= \gnewline{}
<type>? \OPERATOR{} <operator> <formalParameterList>

<оператор>:=~ "
<<P001028> <бинарный оператор>
<<<P001029>> "
<<<P001030><[] "

<binaryOperator> ::= <multiplicativeOperator>
<<p001031> <additive-оператор>
<<P001032> <shiftOperator>
<<P001033> <relationOperator>
<<<P001034>== "
<<<P001035> <bitwiseOperator>
<<P000447>

<<P000737>%
Декларация оператора идентифицируется с помощью встроенного идентификатора.
(<<P000496>)
<<<P001036>.

<<<P000738>%
Для операторов, определяемых пользователем, допускаются следующие имена:
<<<P001037><<,
\lit{>>,
<<<P001039><<==
<<<P001040>{>==
<<<P001041 >>{=======================================================================================================================================================================================================================================================
<<<P001042>{-}
<<<P001043>{+}
<<<P001044>{/}
<<<P001045><<<P001046>/>
\lit{*},
<<<P001048><<<P001279>>
<<<P001049>>
<<<P001050><<<P001280>>
<<<P001051><<<P001281>>
<<<P001052><<<P001053>>
<<<P001054><<<P001055>>,
<<<P001056><<<P001057>>
\lit{[=],
<<<<P001059>{[]],
<<<P001060><<<P001061>>.

<<P000739>%
Ошибка времени компиляции, если оператор объявлен пользователем
<<<P001062>{[]=} не является 2.
Это ошибка времени компиляции, если активность оператора, объявленного пользователем.
с одним из имен:
<<<P001063><<,
\lit{>>,
<<<P001065><<==
<<<P001066>{>==>
<\lit{======================================================================================================================================================================================================================================================
<<<P001068>{-}
\lit{+}
\lit{<<P001071>>>/}
<<<P001072>{/}
\lit{*},
<<<P001074><<<P001282>>
<<<P001075>{ |}}
<<<P001076><<<P001283>>
<<<P001077><<<P001284>>
<<<P001078><<<P001079>>
<<<P001080><<<P001081>>
\lit{\gtgt,
\lit{[]]
Это не 1.
Ошибка времени компиляции, если оператор объявлен пользователем
<<<P001085>{-}
Это не 0 или 1.

<<<P001086>{%>
<<<P001087>{-} Оператор уникален
Допускаются две перегруженные версии.
Если у оператора нет аргументов, он обозначает унарный минус.
Если у него есть аргумент, он обозначает двоичное вычитание.%
?

<<<P000740>%
Имя унарного оператора \lit{-} — <<<P000574>.

<<<P001089>{%>
Это устройство позволяет различать два метода.
для целей метода поиска, переопределения и отражения.%
?

<<<P000741>%
Ошибка времени компиляции, если оператор объявлен пользователем
<<<P001090><<<P001091>>
Это не 0.

<<<P000742>%
Ошибка времени компиляции для объявления необязательного параметра в операторе.

<<<P000743>%
Ошибка времени компиляции, если оператор объявлен пользователем \lit{[=]
Объявляет тип возврата, отличный от <<<P001093>.

<<<P001094>{%>
Если для оператора, объявленного пользователем, не указан тип возврата
<<<P001095>{[=],
Тип возврата - \VOID{} (<<P000497>)%
?

<<<P001097>{%>
Тип возврата <<<P001098>{} Потому что
заявление о возврате в реализации оператора
\lit{[=]
Не возвращает объект.
Рассмотрим оценку выражения <<<P000209> в форме
<<P000575>,
и предположить, что оценка <<<P000213> Дает объекту <<P000214>.
<<P000215> После этого будет оцениваться <<<P000216>,
и даже если исполняемый орган оператора
\lit{[=]
завершается с помощью объекта <<P000217>,
То есть, если <<<P000218> Возвращается, его просто игнорируют.
Обоснование такого поведения заключается в том, что
Задания должны быть гарантированы для оценки назначенному объекту.%
?


<<<P001101> Метод <<<P000576>>
<<P000798>

<<<P000744>%
Метод <<<P000577> Имплицитно упоминается во время исполнения
В ситуациях, когда один или несколько членов поиска терпят неудачу
(<<P000498>,
<<P000499>,
<<P000500>.

<<<P001102>{%>
Мы можем думать о P000578. как резервная копия
который вступает в силу при обращении члена <<<P000219> Попытка,но нет члена, названного <<<P000220>,
или существует,
Но данный призыв имеет форму списка аргументов.
не соответствует декларации <<P000221>
(проведение меньшего количества позиционных аргументов, чем требуется, или больше, чем поддерживается)
или передавать именные аргументы с именами, не объявленными <<<P000222>.
% Следующее предложение охватывает как функциональные объекты, так и экземпляры
% класс с методом, названным <<<P001103>, потому что
% ошибка времени компиляции <<<P001104> с неправильно оформленным аргументом
Список %, если тип не является \DYNAMIC{} или \FUNCTION.
Это может произойти только для обычного вызова метода.
когда приемник имеет статический тип <<<P001107>,
или для выполнения функции, когда
Вызванная функция имеет статический тип <<<P001108>{} или <<<P001109>.
%
Метод <<<P000579> могут быть использованы и другими способами, например,
Его можно назвать как любой другой метод.
На него можно ссылаться из <<<P000580> экспедитор,
Как описано ниже.%
?

<<<P000745>%
Мы говорим, что класс <<<P000223> \Index{имеет нетривиальный \code{noSuchMethod}}
Если <<P000224> имеет конкретный член, названный <<P000582>
который отличается от заявленного во встроенном классе <<<P000583>.

<<<P001111>{%>
Обратите внимание, что это должен быть метод, который принимает один позиционный аргумент.
Чтобы правильно переопределить <<P000584> в <<P000585>.
Например, он может иметь подпись.
<<P000586> или
<<P000587>,
но не
<<P000588>.
Это означает, что ситуация, при которой <<P000589> на него ссылаются
(явно или неявно)
С одним реальным аргументом нельзя не согласиться по той причине, что
«Такого метода нет»
Чтобы мы вошли в бесконечный цикл
Пытаясь вызвать <<P000590>.
Однако можно столкнуться с динамической ошибкой <<<P001112>>
Во время ссылки на <<P000591>
Поскольку фактический аргумент не удовлетворяет проверке типа,
Но эта ситуация приведет к динамическим ошибкам типа.
Вместо того, чтобы повторять попытку <<P000592>
(<<P000501>).
Вот пример, когда ошибка динамического типа возникает из-за
Предпринята попытка пройти <<<P000593>
где принимается только нулевой объект: %
?

<<P000002>

<<<P000746>%
<<<P000817>%
<<<P000225> Будьте конкретным классом,
<<<P000226> быть библиотекой, содержащей декларацию <<<P000227>,
<<<P000228>> Будь именем.
Тогда $m$ \Index{noSuchMethod forwarded} в $C$ если {f}
Одно из следующего верно:

<<P000435>
<<P001113> <<<P001114> Запрошено в программе: }
<<P000231> имеет нетривиальный <<<P000594>,
Интерфейс <<P000232> содержит подпись члена <<P000233> <<<P000234>,
<<<P000235> Нет конкретного члена, названного <<P000236> Доступно для <<P000237>
который правильно перекрывает <<P000238>
<<<P001115>{%>
(то есть, ни один член не назван <<P000239>) Объявляется или наследуется по адресу <<<P000240>,
или наследуется, но не имеет требуемой подписи.
}
В этом случае мы также говорим, что <<<P000241> пересылается noSuchMethod.
<<<P001116>
<<<P001117> Принуждение к конфиденциальности:
Существует прямой или косвенный суперинтерфейс
<<<P000242> <<<P000243> Объявляется в библиотеке <<P000244> Отличие от <<<P000245>,
Интерфейс <<<P000246> содержит подпись члена <<P000247> <<<P000248>>
<<P000249> Это частное имя,
Без суперкласса <<P000250> иметь
Конкретный член, названный <<P000251> Доступно для <<<P000252>
который правильно перекрывает <<P000253>.
В этом случае мы также говорим, что <<<P000254> пересылается noSuchMethod.
<<P000448>

<<<P000747>%
Для конкретного класса <<<P000255>, a
<<P000540>
неявно индуцируется для каждой подписи
Который не является ниспосланным.

<<<P000748>%
Ошибка времени компиляции, если имя <<P000256> NoSuchMethod скачать
в конкретном классе <<<P000257>,и суперкласса <<<P000258> имеет доступную конкретную декларацию <<<P000259>
который не является форвардом.

<<<P000749>%
% % TODO(Eernst): Мы привыкли говорить «интерфейс <<<P000260> или <<<P000261>»; но мы
% необходимо изменить «доступный» на схему, аналогичную слиянию имен
%, чтобы можно было точно сказать.
Экспедитор noSuchMethod является конкретным членом <<<P000262>
с подписью, взятой из интерфейса <<<P000263>,
и с одинаковым значением по умолчанию для каждого дополнительного параметра.
Это может быть вызвано в обычном призыве и в сверхвызове.
и когда <<P000264> Это метод, который можно клозировать.
(<<P000502>)
Использование извлечения имущества
(<<P000503>).

<<<P001118>{%>
Ошибка, связанная с неявно вызванным экспедитором
Это отменяет написанную человеком декларацию.
Это может произойти только в том случае, если конкретное заявление не
правильно отменять <<P000265>. Рассмотрим следующий пример: %
?

<<P000003>

<<<<P001119>{%>
В этом примере,
Реализация с подписью <<P000595> Унаследовано от <<<P000596>,
и суперинтерфейс <<<P000597> объявлять
Подпись <<P000598>.
Это ошибка времени компиляции, потому что <<P000599> не иметь
Реализация метода с подписью <<P000600>.
Мы не хотим косвенно побуждать
<<<P000601> экспедитор с подписью <<P000602>
Потому что он будет перевешивать <<P000603>,
Это, вероятно, будет очень запутанным для разработчиков.
%
В частности, это вызов, как <<<P000604>
призывать <<P000605>,
Хотя это призыв, который является правильным для
Объявление <<<P000606> в <<<P000607>.
%
Поэтому мы требуем от разработчиков четкого разрешения конфликта.
когда неявно индуцируется <<<P000608> экспедитор
будет отменять явно объявленное унаследованное осуществление.
%
Однако это не проблема,
Пусть будет <<<P000609> отклонить
<<<P000610> экспедитор,
В такой ситуации ошибки нет.%
?

<<<P001120>{%>
Это означает, что <<<P000611> У экспедитора такая же
свойства как явно заявленный конкретный член,
За исключением, конечно, того, что <<<P000612> экспедитор
не препятствует индуцированию себя или другого экспедитора noSuchMethod.
Мы не указываем тело <<<P000613> экспедитор,
Но он будет ссылаться на <<<P000614>,
и динамическая семантика ее выполнения указана ниже.%
?

<<<P001121>{%>
В начале этого раздела мы упоминали о неявных призывах.
<<<P000615>> может произойти только
с приемником статического типа \DYNAMIC{}
или функция статического типа \DYNAMIC{} или <<P001124>>>.
<<<P000616> экспедитор,
<<P000617> Также может быть использован
на приемнике, статический тип которого не является <<<P001125>.
Аналогичной ситуации для функций не существует.
потому что невозможно вызвать <<<P000618> экспедитор
в класс функционального объекта.%
?

<<<P001126>{%>
Для конкретного класса <<<P000266>,
Мы можем думать о нетривиальном [P000619].
(объявлено или унаследовано <<P000267>)
в качестве просьбы о "автоматическом осуществлении" всех невыполненных членов
в интерфейсе <<<P000268> в качестве экспедиторов <<<P000620>.
Кроме того, имеется неявная просьба о
Автоматическая реализация всех нереализованных
Недоступные представители любого конкретного класса,
Существует ли нетривиальный P000621?
Обратите внимание, что последний не может быть написан явно в Дарте.
потому что их имена недоступны;
Но язык все еще может указать, что они индуцированы неявно.
потому что компиляторы контролируют обработку личных имен.
?

<<<P000750>%
<<<P000818>%
Для динамической семантики,
Предположим, что класс <<<P000269> имеет неявно индуцированный
<<P000622> экспедитор по имени <<<P000270>,
Формальные параметры типа
<<P000623>,
Формальные позиционные параметры
<<<P000819>%
<<P000624>\commentary{(некоторые из которых могут быть необязательными, когда <<P000273>)},
именованные формальные параметры с именами
<<P000625>
<<<P001128>{(с приведенными выше значениями по умолчанию)}.

<<<P001129>{%>
Для этого не нужно различать
подпись, которая имеет необязательные позиционные параметры и
подпись, в которой указаны параметры,
потому что первое покрывается <<<P000275>.%
?

<<<P000751>%
Исполнение тела <<<P000276> создавать
экземпляр <<<P001130>{im} предопределенного класса <<P000626>
такой, что:

<<P000436>
<<P001131> <<P000627> Показатель <<<P001132> f <<<P000278> Это метод.
<<P001133> <<P000628> Показатель <<<P001134> Если f <<<P000280> является геттером.
<<P001135> <<P000629> Показатель <<<P001136>{} f <<<P000282> Это сеттер.
<<P001137> <<P000630> Оценивает символ <<<P000631>.
<<P001138> <<P000632> Оценка в неизменяемый список
чей динамический тип <<P000633>,
содержат те же объекты, что и список, полученный в результате оценки
<<P000634>.
<<P001139> <<P000635> Оценка на неизменяемую карту
чей динамический тип <<P000636>,
с теми же ключами и значениями, что и карта, полученная в результате оценки

\code{<Symbol, Object><<P001285>>>$\#x_1$: <<P000288>: <<<P000289><<<P001286>>.
<<P001141> <<P000637> Оценить неизменяемый список
чей динамический тип <<P000638>,
содержит те же объекты, что и список, полученный в результате оценки
<<P000639>.
<<P000449>

<<<P000752>%
Далее <<<P000640> Называется <<<P000292> Как реальный аргумент,
и полученный оттуда результат возвращается исполнением <<<P000293>.

<<<P001142>{%>
Это обычный метод вызова <<<P000641>
(<<P000504>).
То есть <<<P000642> экспедитор в классе <<<P000294> может вызвать
Реализация <<<P000643> Об этом говорится в
Подкласс <<<P000295>.

Проверка динамического типа (DVD) - Проверка динамического типа (DVD) - Проверка динамического типа (DVD) - Проверка динамического типа (DVD) - Проверка динамического типа (DVD) - Проверка динамического типа (DVD) - Проверка динамического типа (DVD) - Проверка динамического типа (DVD) - Проверка динамического типа (DVD) - Проверка динамического типа (DVD) - Проверка динамического типа (DVD)
выполняется таким же образом, как и для вызова
Явно заявленный метод.
В частности, фактический аргумент перешел к ковариантному параметру.
Они будут проверяться динамически.

Кроме того, как и другие обычные методы,
Ошибка динамического типа, если результат возвращается
<<P000644> Форвард имеет тип, который не является подтипом.
Тип возврата экспедитора.

Один особый случай, о котором следует знать, это когда экспедитор оторван.
А затем ссылается на фактический список аргументов, который не соответствует
Официальный список параметров.
В этой ситуации мы получим призыв
<<P000645>
Вместо «P000646» в оригинальном приемнике,
потому что это вызов объекта функции
(и они не перекрывают <<<P000647>):%
?

<<P000004>


<<P001143> Оператор «==» и Примитивное равенство
<<P000799>

<<<P000753>%
Оператор <<<P001144>{==} используется неявно в определенных ситуациях,
В частности, постоянные выражения
(<<P000505>)
Это создает ограничения для данного оператора.
Аналогичная ситуация с геттером <<<P000648>.
Для того, чтобы определить эти ограничения, как только мы вводим понятие
<<P000541>.

<<<p001145>{%>
Известно, что некоторые постоянные выражения имеют значение
Равенство которых примитивно.
Это полезно знать, так как позволяет
Значение выражения равенства
и значение вызовов <<<P000649>
вычисляется во время компиляции.
В частности, его можно использовать для строительства
постоянные сборы во время компиляции,
Его можно использовать для проверки
Все элементы в постоянном наборе различны,
и все ключи в постоянной карте различны.%
?

<<P000437>
<<P001146> Нулевой объект имеет примитивное равенство.
(<<P000506>).
<<P001147> Каждый экземпляр типа <<<P000650>, <<P000651> и <<P000652>
примитивное равенство.<<<P001148> Каждый экземпляр типа <<<<P000653>
который был первоначально получен путем оценки буквального символа или
постоянное обращение к конструктору <<<P000654> класс
примитивное равенство.
<<P001149>
%% TODO (Eernst): С Dart 2.15 мы должны упомянуть «Список» здесь,
% %, если «постоянный тип» не включает этот случай.
Каждый экземпляр типа <<<P000655>
которая была первоначально получена путем оценки постоянного типа
(<<P000507>)
примитивное равенство.
<<P001150>
<<<P000297> быть объектом, полученным путем оценки постоянного списка
(<<P000508>),
Постоянная карта буквально
(<<P000509>), или
постоянный набор буквальных
(<<P000510>),
<<<P000298> примитивное равенство.
<<P001151>
% % TODO (Eernst): С Dart 2.15 рассмотрите «f<int>».
Объект функции, полученный путем клосуризации функции
Статический метод или функция верхнего уровня
(<<P000511>)
как значение постоянного выражения
примитивное равенство.
<<P001152>
Пример <<<P000299> примитивное равенство
Если динамический тип <<<P000300> класс <<<P000301>,
<<P000302> примитивное равенство.
<<P001153>
Класс <<<P000656> примитивное равенство.
<<P001154>
Класс <<<P000303> примитивное равенство
если он не имеет реализации оператора <<<P001155>
который перевешивает унаследованный от <<<P000657>,
и не имеет реализации геттера <<<P000658>
Это отменяет унаследованное от P000659.
<<P000450>

<<<P000754>%
Когда мы говорим, что данный случай или класс
<<<P001156>{не имеет примитивного равенства}{%
Примитивное равенство! не имеет,
Это означает, что это неправда, что упомянутый экземпляр или класс
примитивное равенство.


<<<P001157> Геттерс.
<<P000800>

<<<P000755>%
Геттеры — это функции (<<P000512>). которые используются
для получения значений свойств объекта.

<<P000438>
<getterSignature> ::= <type>? \GET{} <идентификатор>
<<P000451>

<<<P000756>%
Если тип возврата не указан, тип возврата геттера <<<P001159>.

<<<P000757>%
Определение Геттера, прикрепленное к <<<P001160>{} Модификатор определяет
Статический геттер.
В противном случае, он определяет пример геттера.
Имя геттера дается идентификатором в определении.

<<<P000758>%
<<<P000542> <<P000304> есть
те экземпляры, которые объявлены <<<P000305>,
явно или косвенно,
и экземпляров, унаследованных <<<P000306> из своего суперкласса.
<<<P000543> <<<P000307> есть
Эти статические геттеры объявлены <<<P000308>.

<<<P001161>{%>
Декларация Геттера может противоречить другим декларациям.
(<<P000513>).
В частности, геттер никогда не может переопределить метод.
и метод никогда не может переопределить геттер или переменную экземпляра.
Правила, когда геттер правильно перевешивает другого члена
Дано в другом месте
(<<P000514>) %
?


<<P001162> Сеттеры.
<<P000801>

<<<P000759>%
Сеттеры — это функции (<<P000515>). которые используются для установки
Значения свойств объекта.

<<P000439>
<setterSignature> ::= <type> \SET{} <identifier> <formalParameterList>
<<P000452>

<<<P001164>{%>
Если тип возврата не указан, тип возврата задаваемого устройства <<<P001165>{}
(<<P000516>)%
?

<<<P000760>%
Определение сеттера, приставленное к <<<P001166>{} Модификатор определяет
Статический сеттер.
В противном случае он определяет набор экземпляров.
Имя сеттера получается путем добавления строки «=» к
идентификатор, указанный в его подписи.

<<<P001167>{%>
Следовательно, имя сеттера никогда не может конфликтовать с, переопределять или переопределяться.
Геттер или метод.%
?

<<<P000761>%
<<<P000544> <<P000309> есть
Настройщики экземпляров, объявленные <<<P000310>
явно или косвенно,
и наборы экземпляров, унаследованные <<<P000311> из своего суперкласса.
<<P000545><<P000312> естьЭти статические установки объявлены <<<P000313>,
либо неявно, либо явно.

<<<P000762>%
Это ошибка времени компиляции, если формальный список параметров множителя
не состоит из одного необходимого формального параметра. <<P000314>.
<<<P001168>{%>
Это можно сделать с помощью грамматики,
Но в этом случае мы должны указать правила оценки.%
?

<<<P000763>%
Это ошибка времени компиляции, если установщик объявляет тип возврата, отличный от <<<P001169>.
Ошибка времени компиляции, если класс имеет
Сеттер по имени <<P000660> Тип аргумента <<<P000316> и
Геттер по имени <<P000317> с возвратным типом <<<P000318>,
<<<P000319> Не может быть назначено <<<P000320>.

<<<P001170>{%>
Правила, когда сеттер правильно перевешивает другого участника
Дано в другом месте
(P000517).
Декларация учредителя может противоречить и другим декларациям.
(<<P000518>) %
?


<<<P001171> Абстрактные участники
<<P000802>

<<<P000764>%
<<<P001172>{абстрактный метод}{метод!абстракт}
(соответственно,
<<<P001173>{abstract getter}{getter!abstract} или
<<<P001174>{abstract setter}{setter!abstract}
является методом экземпляра, геттера или сеттера, который не объявлен <<<P001175>
и не обеспечивает реализацию.
A <<<P001176>{конкретно} Метод {метод!конкретный}
(соответственно,
<<<P001177>{конкретно} Геттер (getter!conrete)
\IndexCustom{concrete setter}{setter!concrete}
Это метод экземпляра, геттера или сеттера, который не является абстрактным.

<<<P001179>{%>
Абстрактные экземпляры полезны из-за их взаимодействия с классами.
Каждый класс Dart создает неявный интерфейс.
Dart не поддерживает явное определение интерфейсов.
Использование абстрактного класса вместо традиционного интерфейса
имеет важные преимущества.
Абстрактный класс может обеспечить реализацию по умолчанию.
Также можно использовать статические методы.
Устранение потребности в классах обслуживания
<<<P000661> или <<P000662>,
цель которой состоит в том, чтобы группировать утилиты, относящиеся к данному типу.
?

<<<P001180>{%>
Призыв абстрактного метода, геттера или сеттера не может произойти.
Поиск (<<P000519>) Никогда не уступать
Абстрактный член в результате.
Один из способов думать об этом — это
Абстрактное заявление члена в подклассе
не отменяет и не отменяет унаследованную реализацию.
Он служит только для указания подписи данного члена, который
каждый конкретный подтип должен иметь реализацию;
То есть он вносит свой вклад в интерфейс класса,
Не для самого класса.%
?

<<<P001181>{%>
Целью абстрактного метода является предоставление декларации.
Для таких целей, как проверка типа и отражение.
В рамках таких формулировок часто используются такие формулировки, которые
ожидаемое смешивание будет обеспечено суперклассом, к которому применяется смесь.%
?

<<<P001182>{%>
Мы хотим определить, объявляется ли конкретный класс абстрактными членами.
Тем не менее, код должен работать следующим образом: %
?

<<P000005>

<<<P001183>{%>
Во время выполнения конкретный метод <<<P000663> Заявлено в <<<P000664>
будут казнены,
И никаких проблем не должно возникнуть.
Поэтому никакой ошибки не должно быть
если соответствующий конкретный член существует в иерархии.
?


<<<P001184> Переменные инстанций
<<P000803>

<<<P000765>%
<<<P001185>{Вариабельные значения}{вариабельные значения!
переменные, декларации которых
немедленно содержатся в классовой декларации
И это не объявлено <<<P001186>.
<<<P000546> <<P000321> есть
Переменные экземпляра, объявленные <<<P000322>
и переменные экземпляра, унаследованные <<<P000323> из своего суперкласса.

<<P000766>%
Это ошибка времени компиляции, если переменная экземпляра объявлена постоянной.

<<<P001187>{%>
Понятие постоянной переменной экземпляра
Это сложно и запутанно для программистов.
Переменная экземпляра предназначена для изменения в каждом экземпляре.Постоянная переменная экземпляра будет иметь одинаковое значение для всех экземпляров.
И как таковой уже является сомнительной идеей.

Язык может интерпретировать заявления с переменной константой как
Геттеры, возвращающие константу.
Однако постоянная переменная экземпляра не может рассматриваться как
Истинная постоянная времени компиляции,
И как бы то ни было, его геттер будет подвержен переопределению.

Учитывая, что стоимость не зависит от конкретного случая,
Лучше использовать статическую переменную.
Геттер экземпляра для него всегда может быть определен вручную, если это необходимо.
?

<<P000767>%
Это возможно для объявления переменной экземпляра
Включает модификатор <<<P001188>{}
(<<P000520>).
Результатом является то, что формальный параметр
соответствующий неявно индуцированный сеттер
считается ковариантным по декларации
(P000521).

<<<P001189>{%>
Модификатор <<<P001190>{} на переменной экземпляра не имеет других эффектов.
В частности, тип возврата неявно индуцированного геттера
может быть преодолена ковариантно без <<<P001191>,
Он никогда не может быть переопределен на супертип или несвязанный тип.
независимо от того, является ли модификатор <<<P001192> Присутствует.%
?


<<P001193> Конструкторы.
<<P000804>

<<P000768>%
<<<P000547> Специальная функция, которая используется
в выражениях создания (<<P000522>) для получения объектов,
Обычно путем их создания или инициализации.
Конструкторы могут быть генеративными (<<P000523>).
Или это могут быть заводы (P000524).

<<P000769>%
<<<P000548> Всегда начинается с
название его непосредственного замкнутого класса,
и необязательно может сопровождаться точкой и идентификатором <<<P001194>.
Ошибка времени компиляции, если имя конструктора
Это не имя конструктора.

<<<P000770>%
The
<<<P001195>{функциональный тип конструктора}{функциональный тип! из конструктора.
<<P000324> является типом функции
класс, который содержит декларацию <<P000325>,
и чьи формальные типы параметров, факультативность и имена названных параметров
соответствует декларации <<<P000326>.

<<<P001196>{%>
Обратите внимание, что тип функции <<<P000327> Конструктор <<P000328> может содержать
переменные типа, объявленные прилагаемым классом <<<P000329>.
В этом случае мы можем применить замену на <<<P000330>, как в
<<P000331>,
где <<P000332> Формальные параметры типа <<P000333>
<<P000334> определяются в данном контексте.
Мы также можем опустить такую замену, когда данный контекст
Объем корпуса <<<P000335>, где <<<P000336> в объеме.%
?

<<<P001197>{%>
Декларация конструктора может противоречить статическим декларациям членов
(<<P000525>)%
?

<<<P001198>{%>
Декларация конструктора не вводит имя в объем.
Если функция выражения вызова
(<<P000526>)
или создание экземпляра
(<<P000527>)
обозначает строителя как
<<P000337>, <<P000665>, <<<P000666>, или \code{\metavar{prefix}.<<<P000341>. <<<P001201>>
разрешение зависит от объема библиотеки для определения класса
(возможно, через префикс импорта).
Классовая декларация проверяется непосредственно
Имеет ли он конструктор под названием <<P000342> соответственно <<P000667>.
Невозможно, чтобы идентификатор непосредственно ссылался на конструктор,
Поскольку конструктор не находится ни в одной области, используемой для разрешения идентификаторов.
?

<<P000771>%
Если для класса <<<P000344> не указан конструктор,
имеет конструктор по умолчанию \code{C():\ \SUPER()\ \{\}},
Если только <<P000345> Это встроенный класс <<P000668>.
% % TODO(Eernst): С нулевой безопасностью добавьте «или <<P000669>».


<<<P001204>> Генеративные конструкторы.
<<P000805>

<<P000772>%
A <<<P001205>{генеративный конструктор}{конструктор!генератор}
Декларация состоит из имени конструктора, формального списка параметров.
(<<P000528>),и либо пункт перенаправления, либо список инициализаторов и факультативный орган.

<<P000440>
<constructorSignature> ::= <constructorName> <formalParameterList>

<constructorName> ::= <typeIdentifier> ('.' <identifier>)?
<<P000453>

<<<<P001206>{%>
Видишь? \synt{declaration} и <<P001208>>>{metodSignature}
введение списка перенаправления или инициализатора и тела.%
?

% % TODO(Eernst): Добавить <<<P001209>{...}> ниже, когда эта команда добавлена.
<<P000773>%
Ошибка времени компиляции возникает, если декларацией генеративного конструктора
Имеет корпус формы <<<P000670>.

<<<<P001210><%
В других декларациях функций,
Этот тип тела подразумевает, что значение <<P000347> возвращается,
Но генеративные конструкторы ничего не возвращают.
?

<<P000774>%
Если официальное заявление о параметрах <<P000348> происходит от
<<<P001211>{поле Формальный Параметр}
Объявлено: <<P000549>.
Термин формы <<<P000671> Содержится в <<P000349>,
<<P000687> <<<P001212>{имя} <<<P000350>.
Ошибка времени компиляции <<<P000688> Не является также именем
переменная экземпляра немедленно закрывающего класса или числа.

<<<P000775>%
Ошибка времени компиляции для инициализации формального параметра
Это происходит в любой функции, которая не является генеративным конструктором.
Кроме того, это ошибка компиляции для инициализации формального параметра.
Происходит в перенаправлении или внешнем конструкторе.
<<<p001213>> В Particuar всегда есть закрытый класс или enum. ?

<<P000776>%
Предположим, что <<P000351> Объявление параметра инициализации, названного <<<P001214>.
<<<P000352> быть типом переменной экземпляра, называемой <<<P000689> в
Немедленно включающий класс или число.
Если <<P000353> имеет аннотацию типа <<<P000354> Заявленный тип <<<P000355> является <<<P000356>.
В противном случае заявленный тип <<<P000357> является <<<P000358>.
Ошибка времени компиляции, если заявленный тип <<P000359>
Не является подтипом <<<P000360>.

<<P000777>%
Инициализация формальностей является исключением из правила, согласно которому
Каждый формальный параметр вводит локальную переменную в
Формальный диапазон параметров (<<P000529>).
Когда формальный список параметров неперенаправляющего генеративного конструктора
содержит любые инициализирующие формы, вводится новый объем,
<<<P001215>{формальный параметр начального диапазона}{%}
Формальный параметр инициализатор.
который является текущим объемом списка инициализаторов конструктора,
и который заключен в область, в которой объявлен конструктор.
Каждая инициализация в официальном списке параметров
вводит конечную локальную переменную в формальную область параметрического инициализатора;
но не в формальный параметр;
любой другой формальный параметр вводит локальную переменную в
Формальный параметр и формальный параметр инициализатор.

<<<<P001216>{%>
Это означает, что формальные параметры, включая инициализацию формальностей,
должны иметь различные названия, и что инициализация формальностей
находятся в пределах списка инициализаторов,
Но они не подходят для корпуса конструктора.
Когда формальный параметр вводит локальную переменную в два объема,
Это все еще одна переменная и, следовательно, одно место хранения.
Тип конструктора определяется по его формальным параметрам.
Включая инициализацию форм.%
?

<<<P000778>%
Инициирование формальностей выполняется во время
Строительные конструкции, описанные ниже.
Выполнение инициализации формально <<P000672>
вызывает переменную экземпляра <<<P000690> непосредственно окружающего класса
для присвоения значения соответствующего фактического параметра,
% Это может произойти из-за неудачного неявного слепка.
Если объект не имеет динамического типа
который не является подтипом заявленного типа переменной экземпляра <<<P001217>,
В этом случае возникает динамическая ошибка.

<<<P001218>{%>Вышеприведенное правило позволяет использовать инициализацию форм в качестве дополнительных параметров: %
?

<<P000006>

<<<<P001219><%
является законным и имеет тот же эффект, что и %
?

<<P000007>

<<<P000779>%
<<<P000550> Это случай, личность которого отличается от
любой ранее выделенный экземпляр своего класса.
Генеративный конструктор всегда работает на свежем примере
сразу же закрывающий класс.

<<<P001220>{%>
Вышеприведенное удерживает, если конструктор фактически запущен, как это происходит на <<<P001221>.
Если конструктор <<P000361> На него ссылается <<<P001222>, <<P000362> не может быть запущен;
Вместо этого можно посмотреть на канонический объект.
См. раздел о создании экземпляра (<<<P000530>).%
?

<<<P000780>%
Если генераторный конструктор <<<P000363> Не является перенаправляющим конструктором
и тело не предоставлено, тогда <<P000364> Имеет пустое тело \code{\{\}}.


<<<P001224> Перенаправление генерирующих конструкторов
<<P000806>

<<<P000781>%
Генератор может быть
<<<P001225>{redirecting}{constructor!redirecting}
В этом случае его единственным действием является вызов другого генеративного конструктора.
У перенаправляющего конструктора нет тела.
Вместо этого он имеет пункт перенаправления, который указывает
К кому обращен призыв,
с какими аргументами.

<<P000441>
<redirection> ::= ':' \THIS{} ('.' <identifier>)? <аргументы>
<<P000454>

\def\ConstMetavar{<<P001229>>>{\CONST?

<<P000782>%
<<<P000820>%
Предположим, что
<<P000673>
- наименование и формальные параметры типа ограждающего класса,
<<<P000821>%
<<P000367> обозначает либо <<<P001231> или ничего,
<<<P000368> - <<<P000369> или <<<P000370> для некоторого идентификатора <<<P000371>,
<<<P000822>%
<<<P000691> - идентификатор.
Рассмотрим декларацию перенаправления генеративного конструктора
одной из форм

<<P001232>
<<P000674>

<<P001233>
<<<P001234><<<P000377> $N$($T_1\ x_1 \ldots,\ T_n\ x_n,\ $\{$T_{n+1}\ x_{n+1} = d_1 \ldots,\ T_{n+k}\ x_{n+k} = d_k$\}):\ $R$;

<<P001235>
где <<P000382> является одной из форм

<<P001236>
<<P000675>

<<P001237>
<<P000676>.

<<P000783>%
The
<<<P001238>{redirectee constructor}{constructor!redirectee}}
Для этого конструктор обозначается
<<P000677> соответственно <<P000678>.
Ошибка времени компиляции, если тип списка статических аргументов
(<<P000531>)
<<<P000679>
не является подходящим совпадением для официального списка параметров перенаправления.

<<<P001239>{%>
Обратите внимание, что в случае, когда не пропущены именованные параметры
Покрывается разрешением <<<P000398> быть нулевым,
В случае, когда <<P000399> является негенерическим классом
Покрывается разрешением <<<P000400> быть нулевым,
В этом случае формальный список параметров типа и фактический список аргументов типа
опущены
(<<P000532>)%
?

<<<P001240>{%>
Нам требуется назначаемый матч, а не более строгий матч подтипа.
потому что конструктор генеративного перенаправления <<<P000401> ссылается на свой редирект <<P000402>
в манере, которая напоминает вызов функции в целом.
Например, <<P000403> Можно ли принять аргумент <<P000680>
и передать выражение <<<P000404>, используя <<<P000681>, такое как <<<P000682>, <<<P000405>,
И это было бы удивительно
Если <<P000406> были подвержены более строгим ограничениям, чем те, которые применялись к
фактические аргументы для функционирования призывов в целом.%
?

<<P000784>%
Перенаправленный генеративный конструктор <<<P000407> <<<P000551>
от перенаправления генеративного конструктора <<P000408> если {f}
<<P000409> является конструктором перенаправления <<<P000410>,
<<<P000411> Конструктор перенаправления <<<P000412>
<<<P000413> Перенаправление доступно из <<<P000414>.
Ошибка времени компиляции при перенаправлении генеративного конструктора
Это перенаправление-доступно от самого себя.

<<<P000785>%
Когда <<<P000415> <<<P001241>,Это ошибка времени компиляции, если перенаправление не является постоянным конструктором.
Более того, когда <<<P000416> <<<P001242>, каждый
<<P000417>,
Должно быть потенциально постоянное выражение (<<P000533>).

<<<P000786>%
% Эта ошибка может произойти из-за неудачного неявного слепка.
Это динамическая ошибка типа, если фактический аргумент принят.
перенаправление генеративного конструктора <<<P000418>
не является подтипом фактического типа (<<P000534>)
соответствующего формального параметра в декларации <<<P000419>.
% Эта ошибка может произойти из-за неудачного неявного слепка.
Ошибка динамического типа, если фактический аргумент принят
перенаправление <<<P000420> перенаправления генеративного конструктора
не является подтипом фактического типа
(<<P000535>)
соответствующего формального параметра в декларации перенаправления.


<<<P001243> Список инициализаторов.
<<P000807>

<<<P000787>%
Список инициализаторов начинается с толстой кишки,
и состоит из запятой отдельного списка отдельных <<<P000552>.

<<<P001244>{%>
Существует три типа инициализаторов.

<<P000442>
<<<P001245>[<<P000421>>>] A <<<P001246>{суперинициализатор} идентифицирует
<<<P001247><<<P001296>--<<<P001297>, то есть
Конкретный конструктор суперкласса.
Исполнение суперинициализаторных причин
Список инициализаторов суперконструктора, который должен быть выполнен.
<<<P001248>[<<P000422>>>] <<<P001249>{фактор инициализации переменной напряжения}
присваивает объект индивидуальной переменной экземпляра.
<<<P001250>[<<P000423>>>] Утверждение.
<<<P000455>%
?

<<P000443>
<initializers> ::= ':' <initializer ListEntry> (',' <инициализатор) ListEntry> *

<инициализатор ListEntry> ::= \SUPER{} <аргументы>
<<P001252> <<<P001253>{} <идентификатор> <аргументы>
<<<P001254> <fieldInitializer>
<<P001255> <утверждение>

<fieldInitializer> ::= <<<P001256><
(<<P001257>>>{} '.')? <идентификатор> '=' <инициализатор Выражение

<initializerExpression> ::= \gnewline{}
<assignableExpression> <assignmentOperator> <expression>
<<P001259> <условное выражение>
<<P001260> <cascade>
<<P001261> <rowExpression>
<<P000456>

% % TODO(Eernst): Когда #54262 будет решен, удалите следующий абзац.
% % Если <выражение> изменяется таким образом, что получается <выражение функции>,
% % следующая ошибка будет просто свойством правила грамматики
%, потому что <initializerExpression> _not_ будет выведено <выражение функции>.
<<P000788>%
В качестве особого правила,
<<<P001262><инициализатор Выражение не может вывести <<<P001263> Выражение.

<<<P001264>{%>
Это решает почти двусмысленность:
В \code{A()<<P001298>>>:\,\,x\,\,=<<P001303>>>\,\,\,\{\,\ldots\,\}},
<<P000683> могут быть инициализированы до пустой записи,
Блок может быть телом конструктора.
Альтернативно, <<<P000684> можно инициализировать в объект функции,
Тогда у строителя не было бы тела.
Было бы известно только, какой случай мы имеем, когда встречаем
(или не встречать)
Семинария в самом конце.
Это считалось нечитаемым.
Следовательно, парсеры могут взять на себя обязательство не анализировать выражение функции.
в этой ситуации.
Обратите внимание, что это все еще возможно для <<<P001267> Выражение} для получения
термин, который содержит выражение функции в качестве подтермина, например,
\code{A()<<P001311>>>:\,\,x\,\,=<<P001316>>>\,(()<<P001318>>>\,\{\,\ldots\,\});}.%
?

<<<P000789>%
Начальник формы <<<P000685> является эквивалентным
инициализатор формы <<P000686>,
Обе формы называются P000553.
Ошибка времени компиляции, если класс
не объявляет переменную экземпляра, названную <<P000428>.Это ошибка времени компиляции, за исключением статического типа <<<P000429>
присваивается заявленному типу <<<P000430>.

